\subsection{How model parameterization changes collapse}
\label{app:theory-size}

The cross-scale losses raise a further question: can a larger model learn
correctness while preserving more correct alternatives? The preceding theorem
already allows every logit to vary independently, so representational capacity
alone cannot answer this question. We give a conditional mechanism: directions
that raise correct logits together can improve correctness with less
amplification of differences among correct modes.

For a smooth parameterization $z=z(\theta)$, write
$J_\theta=\partial z/\partial\theta$ and
$\mathsf K_\theta=J_\theta J_\theta^\top$. The chain rule extends
Lemmas~\ref{lem:grpo-mean} and~\ref{lem:maxrl-mean} to Euclidean parameter
gradient flow:
\begin{equation}
 \dot z=c_G(P)\mathsf K_\theta\nabla_zP,
 \qquad
 \frac{d}{dt}\log\frac{q_c}{q_d}
 =c_G(P)(\mathbf e_c-\mathbf e_d)^\top\mathsf K_\theta\nabla_zP,
 \label{eq:parameterized-collapse}
\end{equation}
where $\mathbf e_c$ is a coordinate basis vector. Thus the logit-gradient geometry,
not the number of columns of $J_\theta$ alone, governs mode amplification.
The independent-logit theorem has $\mathsf K_\theta=I$.

To isolate shared correctness learning, let $v=\1_{\mathcal C}$ be the
vector that is one on correct categories and zero on incorrect categories,
and consider the constant geometry
\begin{equation}
 \mathsf K_\lambda=I+\lambda vv^\top,\qquad \lambda\ge0.
 \label{eq:shared-correctness-kernel}
\end{equation}
The added direction moves all correct logits equally, so it improves their
mass without directly changing their ratios. For example,
$z=u+\alpha\sum_{j=1}^s h_jv$, with trainable $u$ and scalar $h_j$, realizes
this geometry with $\lambda=s\alpha^2$. More such coordinates strengthen the
shared update if $\alpha$ is fixed; normalizing $\alpha=1/\sqrt{s}$ removes
that dependence. These coordinates are redundant for expressivity and assume
perfect correctness alignment. They describe an optimization mechanism, not
a feature that an arbitrary larger language model is guaranteed to possess.

\begin{theorem}[Shared correctness learning preserves more sampled modes at matched accuracy]
\label{thm:shared-correctness}
Compare the flows~\eqref{eq:parameterized-collapse} with
$\mathsf K_\theta=\mathsf K_\lambda$, the same finite initial logits, and any
of the three functions $c_G$ above. Fix a target correct mass
$P_\star\in(P(0),1)$. At the first time each flow reaches $P_\star$,
$\texttt{distinct@}K$ is nondecreasing in $\lambda$ for every integer $K\ge1$,
while $\texttt{pass@}K$ is identical. The increase is strict for $K\ge2$ when
$q(0)$ is nonuniform. In addition, every correct mode satisfies
\begin{equation}
 q_c^{(\lambda)}(P_\star)
 \ge q_c(0)\exp\!\left(-\frac{L}{\lambda+1/m+1/n}\right),
 \quad
 L=\log\frac{P_\star}{1-P_\star}-\log\frac{P(0)}{1-P(0)}.
 \label{eq:shared-correctness-bound}
\end{equation}
Here $m,n$ count output categories, not model parameters. For every fixed
finite $\lambda$, a unique initial maximum still implies eventual
single-mode collapse as $t\to\infty$.
\end{theorem}

\begin{proof}
Let $r_a=p_a/(1-P)$ for $a\in\mathcal N$, and write
$S_q=\sum_cq_c^2$, $S_r=\sum_ar_a^2$. In the effective time
$d\tau=c_G(P)P(1-P)\,dt$, the logit velocities are
$dz_c/d\tau=q_c+\lambda$ and $dz_a/d\tau=-r_a$. Hence
\begin{equation}
 \frac{dq_c}{d\tau}=q_c(q_c-S_q),\qquad
 \frac{dr_a}{d\tau}=-r_a(r_a-S_r),\qquad
 \frac{d}{d\tau}\log\frac{P}{1-P}=S_q+S_r+\lambda.
 \label{eq:shared-correctness-clock}
\end{equation}
The paths $q(\tau),r(\tau)$ therefore do not depend on $\lambda$.
Also $\dot P=c_G(P)P^2(1-P)^2(S_q+S_r+\lambda)>0$; the same interior-limit
argument as in Theorem~\ref{thm:grpo-collapse} gives $P\to1$.
At $P_\star$, the effective time $\tau_\lambda$ uniquely solves
\[
 L=\lambda\tau_\lambda+
   \int_0^{\tau_\lambda}(S_q(u)+S_r(u))\,du.
\]
The right-hand side increases strictly with both its positive time argument
and $\lambda$, so $\tau_\lambda$ decreases strictly with $\lambda$.
Since $S_q\ge1/m$ and $S_r\ge1/n$,
$\tau_\lambda\le L/(\lambda+1/m+1/n)$. Integrating
$d\log q_c/d\tau=q_c-S_q\ge-1$ over this bounded interval gives the probability
floor in Equation~\eqref{eq:shared-correctness-bound}.

At fixed $P_\star$, define
$\mathcal D_K(q)=\sum_c[1-(1-P_\star q_c)^K]$, the expected number of correct modes in
$K$ independent draws. For
$a(u)=KP_\star(1-P_\star u)^{K-1}$, the replicator equation gives
\[
 \frac{d\mathcal D_K}{d\tau}
 =\frac12\sum_{c,d}q_cq_d(q_c-q_d)\bigl(a(q_c)-a(q_d)\bigr)\le0.
\]
The inequality holds because $a$ is nonincreasing; it is strict for $K\ge2$
and nonuniform $q$. Thus the smaller $\tau_\lambda$ gives larger
$\texttt{distinct@}K$, whereas
$\texttt{pass@}K=1-(1-P_\star)^K$ is unchanged. Finally,
$S_q+S_r+\lambda\le2+\lambda$ and $P\to1$ imply $\tau\to\infty$.
The unchanged correct-mode replicator flow then gives eventual collapse
under a unique initial maximum.
\end{proof}

\paragraph{Same accuracy, different remaining breadth.}
Start from \hyperref[ex:theory-running]{the running example} and stop each idealized flow when
$P$ first reaches $0.90$. Both policies then have
$\texttt{pass@8}=1-10^{-8}$. Numerical integration of
Equation~\eqref{eq:shared-correctness-clock} gives the following illustrative
values; the probability floors are calculated separately from
Equation~\eqref{eq:shared-correctness-bound} with $L=\log9$.
\begin{center}
\begin{tabular}{@{}rrrr@{}}
\toprule
$\lambda$ & Conditional mass $q_3$ & Guaranteed $q_3$ floor & $\texttt{distinct@8}$ \\
\midrule
$0$ & $0.05255$ & $0.01924$ & $2.125$ \\
$4$ & $0.08582$ & $0.06623$ & $2.374$ \\
\bottomrule
\end{tabular}
\end{center}
The trajectory values are rounded approximations; the displayed floors are
rounded down. The shared direction preserves more of the initially rare mode
at the same correctness, but its mass still falls below its initial $0.10$.
A better finite-accuracy trajectory can therefore coexist with eventual loss,
as the next corollary makes precise.

\begin{corollary}[Collapse rate as correctness approaches one]
\label{cor:collapse-error-rate}
Under Theorem~\ref{thm:shared-correctness}, fix finite $\lambda\ge0$
and suppose the initial correct-mode maximum is unique. For every other
correct mode $c$,
\begin{equation}
 \lim_{P\uparrow1}\frac{\log q_c^{(\lambda)}(P)}{\log(1-P)}
 =\frac{1}{\lambda+1+1/n}.
 \label{eq:collapse-error-exponent}
\end{equation}
Thus $q_c^{(\lambda)}(P)=(1-P)^{1/(\lambda+1+1/n)+o(1)}$:
larger fixed $\lambda$ reduces the decay exponent, but every minority mode
still vanishes. The exponent describes the limiting ratio of logarithms;
finite-accuracy probabilities also depend on the initial distribution.
\end{corollary}

\begin{proof}
Theorem~\ref{thm:shared-correctness} gives $q\to\mathbf e_{c^\star}$,
so $S_q\to1$. The incorrect conditional flow in
Equation~\eqref{eq:shared-correctness-clock} preserves ordering and makes
$r_{\min}$ nondecreasing. For $h=\log(r_{\max}/r_{\min})$,
\[
 h'=-(r_{\max}-r_{\min})
 =-r_{\min}(e^h-1)\le-r_{\min}(0)h.
\]
Hence $r\to(1/n,\ldots,1/n)$ and $S_r\to1/n$, including trivially $n=1$.
Integrating $d\log q_c/d\tau=q_c-S_q$ and the correctness-clock equation
and taking time averages yields
\[
 \frac{\log q_c(\tau)}{\tau}\to-1,
 \qquad
 \frac{\log(P(\tau)/(1-P(\tau)))}{\tau}
 \to\lambda+1+1/n.
\]
Since $\log(1-P)=\log P-\log(P/(1-P))$ and $\log P\to0$, their ratio proves
Equation~\eqref{eq:collapse-error-exponent}.
\end{proof}

\paragraph{A smaller exponent still permits extinction.}
With one incorrect category, the minority-mode exponents are $1/2$ at
$\lambda=0$ and $1/6$ at $\lambda=4$. Both are positive, so both limits
lose minority modes as $P\uparrow1$. These are asymptotic logarithmic slopes;
they do not replace the finite-accuracy values in the preceding table.

\paragraph{Correctness alone does not imply collapse.}
The independent-logit component matters. With purely shared geometry
$\mathsf K=\lambda vv^\top$, $\lambda>0$, all correct logits move equally,
so $q(t)=q(0)$ while
$\dot P=\lambda c_G(P)P^2(1-P)^2>0$ and $P\to1$.
For $\mathsf K=aI+bvv^\top$ with fixed $a>0$ and $b\ge0$, time rescaling
instead recovers Theorem~\ref{thm:shared-correctness} with $\lambda=b/a$.
Thus persistent mode-specific amplification, not increasing correctness by
itself, drives this collapse result.

\paragraph{Why parameter count alone is insufficient.}
For an integer $h\ge1$, replace each independent logit by
$z_a=h^{-\gamma}\sum_{j=1}^h\theta_{aj}$, initializing
$\theta_{aj}(0)=h^{\gamma-1}z_a(0)$. Every added parameter affects the logits,
and every model represents the same categorical policies, but
$\mathsf K=h^{1-2\gamma}I$. Its policy follows exactly the original collapse
trajectory with time multiplied by $h^{1-2\gamma}$. Increasing the parameter
count therefore makes collapse faster, unchanged, or slower for
$\gamma=0,1/2,1$, respectively. For example, using $h=4$ multiplies the speed by $4$, $1$, or $1/4$
under those three normalizations, despite quadrupling the parameter count
in every case. A size claim therefore requires assumptions about how
additional parameters change the update geometry.

\paragraph{What this explains and what remains untested.}
Theorem~\ref{thm:shared-correctness} connects a specified kind of shared
learning to the paper's sampling metrics at matched accuracy. It also
clarifies why a slower loss can coexist with eventual collapse, leaving a
role for replay. It does not identify the cause of the empirical cross-scale
pattern: those runs start from different pretrained distributions, use
different learning-rate schedules, and compare equal training passes rather
than matched accuracy. Falcon also changes model family. Neither the assumed
geometry nor its dependence on size is measured here. Establishing that link
would require controlled within-family scaling and measurements of shared
versus mode-specific updates. The result retains prompt isolation, exact
mean updates, and zero reference KL; it does not cover shared-prompt
interference, adaptive optimizer effects, or finite-step stochastic training.

